\section{Conclusion and Future Work}
\label{sec:conclusion}

US early learning standards are public policy that software cannot read. I have
shown that the obstacle is not text recovery but hierarchy recovery, and that an
LLM can recover the hierarchy by classifying each element's nesting position
against a depth map inferred per document, rather than by recognizing label
vocabulary that no two states share. Across four development states and two held
out entirely, detector recall at the annotated subset tier is 1.000 at every
level, a manual audit of all 297 in-scope detections there finds one
hallucinated element, and parsing resolves those
elements into a single canonical schema with no identifier collisions. Ablating
the depth map degrades exactly the levels whose identity is positional and
leaves the others untouched; a rule-based extractor matches the method at a
document's top division and collapses beneath it. The system carries no
per-state code path, which is what allows a new state to be added without new
code; the prompts' worked examples are the one place the corpus's own documents
appear, and they are itemized rather than left implicit.

Quality at full-document scale, which an earlier draft named as the measurement
this work most conspicuously lacked, is now reported: on Kentucky's trimmed tier
all 277 elements of a verified reference set are recovered in each of three
independent runs, with no hallucination, and parser field accuracy is $0.995$.
What that leaves is a bounded and stateable gap rather than an absence --- the
reference set was verified by inspecting what the system emitted, so it
establishes precision and not recall, and a cold read of a random sample of
pages is the cheap way to close it. Beyond that, the corpus is six states of
more than fifty, and the obvious direction is breadth: each state added tests
generalization again, and enough of them would constitute the national
machine-readable dataset that does not currently exist.

Two further directions follow from the artifact rather than the method.
Exposing the normalized store through a tool-callable interface would let
external applications query standards directly, which is the form curriculum
and planning tools need them in. And the per-element verification state invites
a targeted-review question I have deliberately not answered: given that
confidence is uninformative (\S\ref{sec:experiments-confidence}), what
\emph{would} identify the records most worth an expert's attention? My
stability measurements suggest one candidate --- disagreement across repeated
runs is cheap to compute and, unlike confidence, is concentrated in exactly the
places where my own defects live.
